\documentclass[preprint,12pt]{elsarticle}

% add figure packages
\usepackage{graphicx}
\graphicspath{{figures/}} % tells latex where to look

% add algorithm packages
\usepackage{algorithm}
\usepackage{algpseudocode}

% add booktabs package for better table formatting
\usepackage{booktabs}

% The amssymb package provides various useful mathematical symbols
\usepackage{amssymb}

% The amsmath package provides various useful equation environments.
\usepackage{amsmath}
% The amsthm package provides extended theorem environments
% \usepackage{amsthm}

% Add packages for URL and hyperref support
\usepackage{url}
\usepackage{hyperref}
\usepackage{cleveref}

% Fix for float placement in figures
\usepackage{float}

%% The lineno package adds line numbers. 
\usepackage{lineno}
% Line numbers only for Introduction to Conclusion
\usepackage{etoolbox}
\AtBeginEnvironment{document}{}
\newcommand{\startmainlinenumbers}{\begin{linenumbers}}
\newcommand{\stopmainlinenumbers}{\end{linenumbers}}

\journal{Automation in Construction}

\begin{document}

\begin{frontmatter}

%% Title, authors and addresses
\title{ResPointNet-LT: A Deep Long-Tailed Learning Framework for Semantic Segmentation of Large-scale Industrial Point Clouds}
% Alternative titles:
% \title{Long-Tailed Learning for Industrial Point Cloud Semantic Segmentation: A ResPointNet++ Based Approach}
% \title{Addressing Class Imbalance in Industrial Point Cloud Segmentation with ResPointNet-LT}
% \title{ResPointNet-LT: Two-Stage Decoupled Training for Long-Tailed Point Cloud Segmentation}

\author[label1]{Chao Yin}
\author[label2]{Qing Han}
\author[label3]{Zhiwei Hou}
\author[label4]{Boyu Wang}
\author[label5]{Shuju Jing}
\author[label1,label3]{Weilong Jing}
\author[label1,label3]{Ji Yang\corref{cor1}}

\cortext[cor1]{Corresponding author. Email: TO BE FILLED}

%% Author affiliation
\affiliation[label1]{organization={Guangzhou Institute of Geography, Guangdong Academy of Sciences},
            addressline={100 Xianlie Middle Road}, 
            city={Guangzhou},
            postcode={510070}, 
            state={Guangdong},
            country={China}}
\affiliation[label2]{organization={School of Geography and Tourism, Hengyang Normal University},
            addressline={No. 36, Xueyuan Road},
            city={Hengyang},
            postcode={421002},
            state={Hunan},
            country={China}}
\affiliation[label3]{organization={Southern Marine Science and Engineering Guangdong Laboratory (Guangzhou)},
            addressline={},
            city={Guangzhou},
            postcode={510070},
            state={Guangdong},
            country={China}}
\affiliation[label4]{organization={Department of Civil and Environmental Engineering, The Hong Kong University of Science and Technology},
            addressline={Clear Water Bay},
            city={Hong Kong},
            postcode={},
            state={},
            country={China}}
\affiliation[label5]{organization={School of Civil, Architectural and Environmental Systems Engineering, Sungkyunkwan University},
            addressline={2066 Seobu-ro, Jangan-gu},
            city={Suwon-si},
            postcode={16419},
            state={Gyeonggi-do},
            country={Republic of Korea}}


%% Abstract
\begin{abstract}
% Importance of PC SemSeg – Class imbalance problem – Previous work limitation – Proposal – Core components – Results – Research meaning
Point cloud semantic segmentation (PCSS) plays a crucial role in understanding complex 3D industrial environments, enabling automated analysis of large-scale infrastructure and manufacturing scenes. While existing deep learning-based PCSS methods have achieved remarkable progress on balanced benchmark datasets, they suffer from significant performance degradation when applied to real-world industrial point clouds that exhibit severe class imbalance. In industrial scenes, critical components such as pumps, valves, and sensors are dramatically under-represented compared to dominant structural elements like pipes and tanks, creating a long-tailed distribution that poses significant challenges for accurate segmentation of minority classes.
To address this critical limitation, we propose ResPointNet-LT, a novel two-stage decoupled training framework that effectively tackles the class imbalance problem in industrial point cloud segmentation. Our approach builds upon the proven ResPointNet++ architecture and introduces a systematic long-tailed learning strategy specifically designed for 3D point clouds. The proposed ResPointNet-LT framework comprises two key stages: (1) representation learning stage that trains the full network on the original imbalanced dataset to learn rich geometric features, and (2) classifier re-training stage that freezes the encoder weights and fine-tunes only the decoder layers using class-balanced sampling to adjust decision boundaries without corrupting learned representations.
We conduct comprehensive experiments on large-scale industrial point cloud datasets to validate the effectiveness of our proposed method. The experimental results demonstrate that ResPointNet-LT significantly outperforms the baseline ResPointNet++ and other state-of-the-art methods, achieving substantial improvements in tail class IoU (XX\% improvement on pump class, XX\% on valve class) while maintaining strong performance on head classes, resulting in an overall mIoU improvement of XX\%. Extensive ablation studies further validate the contribution of each component in our framework. This study contributes to practical deployment of deep learning techniques in real-world industrial applications such as digital twin construction, infrastructure maintenance, and automated facility management.
\end{abstract}

%%Graphical abstract
\begin{graphicalabstract}
\begin{figure}[H]
\centering
\includegraphics[width=\textwidth]{figures/graphical_abstract.png}
\end{figure}
\end{graphicalabstract}

%%Research highlights
\begin{highlights}
\item Proposes ResPointNet-LT, a novel two-stage decoupled training framework for addressing class imbalance in industrial point cloud segmentation.
\item Identifies and systematically analyzes the severe long-tailed distribution problem in real-world industrial point cloud datasets.
\item Develops a systematic approach to decouple representation learning from classifier training, specifically adapted for hierarchical point-based networks.
\item Achieves substantial improvements in tail class performance (XX\% IoU improvement on minority classes) while maintaining head class accuracy.
\item Demonstrates superior performance on large-scale industrial datasets, advancing practical deployment of deep learning in real-world industrial applications.
\end{highlights}

%% Keywords
\begin{keyword}
Point cloud semantic segmentation \sep Long-tailed learning \sep Class imbalance \sep Industrial point clouds \sep Decoupled training \sep Deep learning
\end{keyword}
\end{frontmatter}

%% main text
%%%%%%%%%%%%%%%%%%%%%%%%%%%%%%%%%%%%%%%%%%
\startmainlinenumbers
%% introduction
\section{Introduction}
With the rapid advancement of 3D hardware devices such as LiDAR sensors, UAVs, and RGB-D cameras, acquiring point clouds of real-world industrial scenes has become increasingly cost-effective and convenient. This technological progress has led to the accumulation of vast amounts of industrial point cloud data from complex facilities, manufacturing plants, and infrastructure systems, presenting both opportunities and challenges for automated analysis and understanding. Concurrently, the remarkable strides made in deep learning have propelled the development of point cloud understanding algorithms, enabling more sophisticated and accurate interpretations of 3D industrial environments. The convergence of these hardware and software advancements has catalyzed a wide range of industrial applications that leverage point clouds to extract valuable information from complex facilities. From digital twin construction~\cite{misc} and predictive maintenance~\cite{misc} to automated facility management~\cite{misc}, quality control~\cite{misc}, and industrial robotics~\cite{misc,misc}, the ability to understand complex industrial 3D scenes from point clouds has become a critical research focus for both academia and industry.

One fundamental task in industrial point cloud scene understanding is point cloud semantic segmentation (PCSS), which empowers machines to interpret and reason about complex industrial environments with fine-grained precision. By accurately segmenting and classifying individual points within an industrial point cloud, PCSS enables machines to identify and understand various components such as pipes, valves, pumps, sensors, and structural elements, unlocking a myriad of applications and paving the way for more intelligent and autonomous industrial systems.

Recently, deep neural networks for point cloud semantic segmentation have achieved tremendous progress, owing to their scalability with large, well-annotated training datasets~\cite{misc,misc,misc}. Pioneering works such as PointNet~\cite{qi2017pointnet} and PointNet++~\cite{qi2017pointnetpp} established the foundation for point-based neural networks, with PointNet++ introducing hierarchical feature learning that captures multi-scale geometric structures. Building upon these seminal works, numerous advanced architectures have been proposed, including RandLA-Net~\cite{hu2020randlanet}, KPConv~\cite{thomas2019}, PointTransformer~\cite{zhao2021pointtransformer}, and PAConv~\cite{xu2021paconv}. Our previous work, ResPointNet++~\cite{respointnet}, further enhanced the PointNet++ architecture with residual connections and improved local feature aggregation, demonstrating superior performance on industrial point cloud segmentation tasks.

However, despite these advances, a critical and underexplored challenge emerges when applying these sophisticated networks to real-world industrial point clouds: severe class imbalance and long-tailed distribution. Unlike carefully curated benchmark datasets that often maintain relatively balanced class distributions, real-world industrial scenes exhibit extreme class imbalance where dominant structural elements (such as pipes, walls, and floors) comprise the vast majority of points, while critical functional components (such as pumps, valves, sensors, and control panels) are dramatically under-represented. This long-tailed distribution poses significant challenges for deep learning models, which tend to be biased toward head classes and perform poorly on tail classes that are often the most critical for industrial applications. For instance, in typical industrial facilities, pipes may constitute 60-70\% of all labeled points, while critical components like pumps and sensors may represent less than 1\% each, creating an imbalance ratio exceeding 100:1.

\begin{figure}[htbp]
    \centering
    \begin{minipage}{0.48\textwidth}
        \centering
        \textbf{Head classes (first 5 classes):}\\
        pipe, wall, floor, structure, platform\\[1ex]
        \textbf{Tail classes (last 5 classes):}\\
        pump, valve, sensor, control\_panel, instrument\\
        (a) Long-tail distribution in Industrial Dataset
    \end{minipage}
    \hfill
    \begin{minipage}{0.48\textwidth}
        \centering
        \textbf{ResPointNet++:} pipe, wall, floor, structure, platform\\
        \textbf{RandLA-Net:} pipe, wall, floor, structure, platform\\
        (b) Baseline models' best performing class types\\[1ex]
        \textbf{ResPointNet++:} pump, valve, sensor, instrument, control\_panel\\
        \textbf{RandLA-Net:} pump, valve, sensor, instrument, control\_panel\\
        (c) Baseline models' worst performing class types
    \end{minipage}
    \caption{(a) Industrial point cloud dataset exhibits severe long-tail distribution with critical components being dramatically under-represented. (b) Existing deep learning models (ResPointNet++ and RandLA-Net) perform well on head classes that dominate the dataset. (c) The same models struggle significantly with tail classes that are often the most critical for industrial applications. This demonstrates the urgent need for specialized techniques to address class imbalance in industrial point cloud segmentation.}
    \label{fig:longtail}
\end{figure}

To address this critical limitation, this paper presents ResPointNet-LT, a novel two-stage decoupled training framework that systematically tackles the class imbalance problem in industrial point cloud segmentation. Our approach builds upon our previously established ResPointNet++ architecture and introduces pioneering long-tailed learning strategies specifically adapted for 3D point clouds. By addressing the long-tailed distribution challenge, our proposed method paves the way for more robust and practical point cloud semantic segmentation, unlocking the potential for widespread deployment in real-world industrial applications.

The core innovation of ResPointNet-LT lies in a systematic decoupling of representation learning from classifier training, specifically adapted for hierarchical point-based networks. First, we identify the fundamental challenge that aggressive re-weighting or re-sampling methods can corrupt the learned geometric representations in point-based networks due to the importance of spatial relationships. Second, we propose a two-stage training strategy: Stage 1 focuses on representation learning where the full ResPointNet++ network is trained on the original imbalanced dataset to learn rich geometric features from all available data; Stage 2 focuses on classifier re-training where we freeze the encoder weights and fine-tune only the decoder layers using class-balanced sampling to adjust decision boundaries without corrupting the learned spatial representations. Third, we introduce enhanced evaluation metrics including per-class IoU analysis and Tail-mIoU to provide comprehensive assessment of model performance on both head and tail classes. Finally, comprehensive experiments on large-scale industrial point cloud datasets demonstrate the effectiveness of our proposed ResPointNet-LT framework compared to strong baselines.

To conclude, our primary contributions are as follows:
\begin{itemize}
    \item \textbf{Problem identification and analysis:} We identify and systematically analyze the severe class imbalance problem in real-world industrial point cloud datasets, demonstrating its significant impact on segmentation performance of critical components.
    \item \textbf{Novel two-stage framework:} We propose ResPointNet-LT, a pioneering two-stage decoupled training framework that systematically addresses class imbalance in hierarchical point-based networks by separating representation learning from classifier training.
    \item \textbf{3D-specific long-tailed learning:} We adapt and extend long-tailed learning techniques from 2D image domain to 3D point clouds, addressing unique challenges such as spatial relationships and geometric feature preservation in point-based networks.
    \item \textbf{Comprehensive evaluation:} We demonstrate that ResPointNet-LT significantly outperforms strong baselines including our previous ResPointNet++ work, achieving substantial improvements in tail class performance while maintaining head class accuracy.
    \item \textbf{Thorough ablation studies:} We perform extensive ablation studies analyzing the contribution of each component, different loss functions, training strategies, and hyperparameter choices in the ResPointNet-LT framework.
\end{itemize}

This paper is organized as follows: \autoref{sec:related_works} reviews deep learning-based semantic segmentation methods for point clouds, with emphasis on fully supervised approaches, and comprehensively surveys long-tailed learning techniques for both 2D images and 3D point clouds. \autoref{sec:methodology} details the proposed ResPointNet-LT framework, including the two-stage decoupled training strategy and evaluation metrics. \autoref{sec:results} presents comprehensive experimental results and ablation studies on industrial point cloud datasets. Finally, \autoref{sec:conclusion} concludes the paper and outlines future research directions.

%% related works
\section{Related works}
\label{sec:related_works}

In this section, we provide a comprehensive review of point cloud semantic segmentation methods with particular emphasis on fully supervised approaches (Section~\ref{sec:fully_supervised}), followed by an in-depth survey of long-tailed learning techniques for both 2D images and 3D point clouds (Section~\ref{sec:long_tailed}). We also discuss specific challenges and existing solutions for industrial point cloud analysis (Section~\ref{sec:industrial_challenges}).

\subsection{PCSS Using Fully Supervised Learning}
\label{sec:fully_supervised}

Fully supervised PCSS methods learn representations from input point clouds using complete, point-wise labeled data. These methods have become dominant in 3D semantic segmentation due to their ability to extract useful information from point clouds without complex feature engineering, their strong capacity to model complex patterns and relationships in large-scale data, and their suitability for analyzing complex 3D environments~\cite{guo2020deep}. Over the past five years, deep learning-based PCSS methods have made tremendous progress: numerous high-quality, well-annotated datasets have been established, and a wealth of sophisticated local aggregation operators have been developed to extract patterns from 3D scenes. Based on the input representation, fully supervised PCSS methods can be broadly divided into three categories:
\begin{enumerate}
    \item \textbf{Point-based methods:} These operate directly on raw point clouds and have achieved impressive results on several benchmark datasets~\cite{qi2017pointnet, hu2020randla, yin2021paconv}. Compared to other types, point-based networks do not require data conversion, making them relatively easy to implement. There is a rich variety of segmentation network designs and local aggregation operators for effective point cloud representation learning.
    \item \textbf{Projection-based methods:} These methods first project the point cloud into 2D images and then apply mature 2D CNNs for segmentation~\cite{qi2016volumetric}. However, the projection process can result in information loss, and selecting appropriate projection angles is practically challenging.
    \item \textbf{Voxel-based methods:} These methods voxelize the point cloud and apply 3D CNNs~\cite{graham2017submanifold}. Voxelization is computationally expensive and memory-intensive, making it less popular for large-scale point cloud segmentation.
\end{enumerate}

Although fully supervised methods can achieve impressive segmentation performance on balanced datasets, they suffer significant performance degradation when applied to real-world datasets with severe class imbalance. The challenge is not the lack of labeled data, but rather the highly skewed distribution where minority classes are dramatically under-represented. This study addresses this critical limitation by developing specialized training strategies that enable fully supervised methods to achieve robust performance across all classes, particularly improving the recognition of under-represented but critical components.

\subsubsection{Point-based Methods}
Point-based methods extract features or embeddings directly from raw points without data conversion, making them the most natural approach for processing unstructured point clouds. Due to the irregular and unordered nature of point clouds, standard 2D CNNs cannot be directly applied. To address this fundamental challenge, PointNet~\cite{qi2017pointnet} was proposed as the pioneering work, enabling feature learning on irregular point clouds via shared multilayer perceptrons (MLPs) and a symmetric aggregation operation (e.g., max-pooling) to obtain permutation-invariant global information. However, PointNet does not explicitly model local context and spatial relationships, which are crucial for complex scene understanding.

To address this limitation, PointNet++~\cite{qi2017pointnet++} was proposed, which extends PointNet by incorporating local context through hierarchical feature learning at multiple scales and locations. PointNet++ introduces the set abstraction layer that recursively applies PointNet on nested partitioning of the input set, enabling the network to learn local features with increasing contextual scales. This hierarchical design has proven fundamental to most subsequent point-based architectures and forms the backbone of our previous work.

Building upon PointNet++, numerous advanced architectures have been developed: DGCNN~\cite{wang2019dynamic} uses dynamic graphs to explicitly learn local geometric structures; PointCNN~\cite{li2018pointcnn} introduces the X-Conv operator that generalizes 2D convolution to point clouds; KPConv~\cite{thomas2019kpconv} proposes kernel point convolution with deformable kernels; RandLA-Net~\cite{hu2020randla} introduces efficient random sampling and local feature aggregation for large-scale point clouds; PAConv~\cite{xu2021paconv} designs position-adaptive convolutions; and PointTransformer~\cite{zhao2021pointtransformer} adapts transformer architecture to point clouds.

Our previous work, ResPointNet++~\cite{respointnet}, significantly enhanced the PointNet++ architecture by integrating deep residual connections, improved local feature aggregation operators, and multi-scale feature fusion mechanisms specifically designed for complex industrial scenes. ResPointNet++ demonstrated superior performance on industrial point cloud segmentation tasks, achieving state-of-the-art results on several benchmark datasets. However, like other fully supervised methods, ResPointNet++ suffers from performance degradation when applied to imbalanced datasets where tail classes are severely under-represented. This limitation motivates our current work to develop ResPointNet-LT, which specifically addresses the class imbalance challenge while building upon the proven ResPointNet++ architecture.

\subsubsection{Projection-based Methods}
Projection-based methods learn features on 2D images rendered from point clouds using mature 2D CNNs. Hang et al.~\cite{su2015mvcnn} proposed Multi-view Convolutional Neural Networks (MVCNNs), which fuse information from multiple views into a compact global shape descriptor via max-pooling. However, max-pooling only preserves the maximum value per view, leading to information loss. Yu et al.~\cite{yu2018multi} introduced bilinear pooling to fuse local convolutional features for a more discriminative global descriptor. Ze and Liwei~\cite{zhang2019relation} presented a Relation Network to join corresponding regions across different views. Other methods~\cite{wang2019dominant, feng2018gvcnn, shi2019pvcnn, liu2019relation} have also been proposed to improve object recognition accuracy. However, projection-based methods are not widely used for point cloud segmentation due to geometric and structural information loss, sensitivity to occlusion, and the challenge of selecting camera views~\cite{qi2016volumetric}.

\subsubsection{Voxel-based Methods}
Voxel-based methods convert point clouds into voxels to leverage 3D CNNs. Maturana and Scherer~\cite{maturana2015voxnet} proposed VoxNet, a volumetric representation-based network for voxel classification. Wu et al.~\cite{wu20153d} introduced 3D ShapeNets, a convolutional deep belief network for learning binary probability distributions of 3D voxels. However, these methods are computationally and memory intensive. To address this, hierarchical structures such as Octree are used. Riegler et al.~\cite{riegler2017octnet} proposed OctNet, which uses unbalanced octrees for efficient memory allocation and computation. Despite these advances, voxel-based methods still suffer from information loss and discretization issues, and selecting suitable grid resolutions requires empirical tuning.

In summary, while fully supervised point-based methods have achieved remarkable success on balanced benchmark datasets, they face significant challenges when applied to real-world industrial scenarios where class imbalance is severe and critical components are under-represented. This motivates the need for specialized training strategies that can leverage the full potential of these proven architectures while addressing the fundamental challenge of long-tailed distributions.

\subsection{Industrial Point Cloud Challenges}
\label{sec:industrial_challenges}

Industrial point cloud analysis presents unique challenges that distinguish it from general indoor or outdoor scene understanding. These challenges stem from the complex nature of industrial facilities, the critical importance of minority components, and the specific requirements for practical deployment.

\textbf{Class Imbalance and Long-Tailed Distributions:} Industrial scenes exhibit severe class imbalance where structural elements like pipes, walls, and platforms dominate the point cloud, often comprising 70-90\% of all points. Critical functional components such as pumps, valves, sensors, and control systems are dramatically under-represented, typically accounting for less than 5\% of the total points. This creates imbalance ratios exceeding 100:1, which is far more severe than typical computer vision datasets.

\textbf{Critical Minority Classes:} Unlike general scene understanding where all classes may have similar importance, industrial applications often prioritize the accurate identification of minority components. Pumps, valves, and sensors are critical for facility operation and maintenance, making their accurate segmentation essential despite their scarcity in the training data.

\textbf{Complex Geometric Structures:} Industrial facilities contain complex geometric arrangements with high levels of occlusion, intricate pipe networks, and densely packed equipment. The spatial relationships between different components are crucial for understanding system functionality.

\textbf{Scale and Density Variations:} Industrial point clouds often span large areas with significant variations in point density. Critical components may be represented by relatively few points compared to large structural elements, creating an additional challenge beyond simple class frequency imbalance.

\subsection{Deep Long-Tailed Learning}
\label{sec:long_tailed}

Long-tailed learning addresses the fundamental challenge of training deep neural networks on datasets with imbalanced class distributions, where a few classes (head classes) dominate the dataset while many classes (tail classes) are severely under-represented. This problem is ubiquitous in real-world applications including computer vision~\cite{zhang2021deep}, natural language processing, and more recently, 3D point cloud analysis~\cite{armeni2016s3dis, dai2017scannet, hackel2017semantic3d}. The class imbalance creates a significant challenge for deep learning models, which tend to be biased toward head classes during training, resulting in poor generalization to tail classes. Recent advances in long-tailed learning have made significant progress in addressing this challenge~\cite{zhang2021deep, zhang2021distribution}. We provide a comprehensive review of long-tailed learning techniques in both 2D images and 3D point clouds, with particular emphasis on methods relevant to our proposed framework.

\subsubsection{2D Image Long-Tailed Learning Methods}
Long-tailed learning methods for 2D images can be systematically categorized into four main families based on their core strategies:

\textbf{Re-weighting Methods} assign different loss weights to classes to focus training on under-represented tail classes. The focal loss~\cite{lin2017focal} addresses class imbalance by dynamically adjusting the cross-entropy loss to down-weight easy examples and focus on hard examples. Class-balanced loss~\cite{cui2019class} introduces effective number of samples to re-weight the loss function. However, re-weighting methods may be ineffective when class imbalance is severe, and aggressive re-weighting can sometimes hurt overall performance.

\textbf{Re-sampling Methods} modify the training data distribution by over-sampling tail classes or under-sample head classes to create more balanced training batches. These include simple random over-sampling, SMOTE~\cite{chawla2002smote}, and more sophisticated sampling strategies~\cite{more2016survey, buda2018systematic}. While effective in some cases, re-sampling methods can introduce noise, overfitting to tail classes, and may not address the fundamental representation learning challenges in imbalanced datasets.

\textbf{Data Augmentation and Mix-up Methods} use sophisticated data augmentation techniques to increase diversity and improve generalization. Input mix-up~\cite{zhang2018mixup} linearly combines input images and their labels, while manifold mix-up~\cite{verma2019manifold} performs mixing in the hidden representation space. These methods are relatively easy to implement and can effectively regularize models while improving tail class recognition.

\textbf{Two-stage and Decoupled Training Methods} represent a breakthrough approach that separates representation learning from classifier training. Kang et al.~\cite{kang2020decoupling} demonstrated that decoupling training can significantly improve long-tailed performance by first training the entire network on imbalanced data to learn good representations, then re-training only the classifier on balanced data. Two-stage training~\cite{cao2019learning} follows a similar philosophy but uses different balancing strategies. These methods have shown remarkable success and form the foundation of our proposed ResPointNet-LT framework.

The success of decoupled training in 2D computer vision motivates its adaptation to 3D point clouds. However, the unique characteristics of point cloud data, including irregular structure, varying density, and the importance of spatial relationships, require careful consideration when applying these techniques to point-based networks.

\subsubsection{3D Point Cloud Long-Tailed Learning Methods}
Long-tailed learning for 3D point clouds remains significantly under-explored compared to its 2D counterparts, with only a few pioneering works addressing this critical challenge~\cite{zhu2019class, chen2022automatic, ding2022domain}. The limited research in this area partly stems from the unique challenges posed by point cloud data structure and the relatively recent emergence of large-scale point cloud datasets.

Zhu et al.~\cite{zhu2019class} proposed a class-balanced multi-head network specifically designed for autonomous driving scenes, addressing the imbalance between common objects (cars, roads) and rare objects (motorcycles, bicycles). Their approach uses multiple classification heads with different sampling strategies. Chen et al.~\cite{chen2022automatic} developed an annotation workflow for synthetic point clouds, intentionally placing more tail-class objects in scenes to create more balanced training data. While innovative, this approach requires manual scene construction and may not capture the natural distribution of real-world environments.

Ding et al.~\cite{ding2022domain} introduced domain adaptation techniques combined with tailed data cuboid mixing to improve segmentation performance on minority classes. Their method focuses on cross-domain transfer learning but does not fundamentally address the within-domain class imbalance problem.

Despite these early efforts, several critical challenges remain unaddressed in 3D point cloud long-tailed learning:

\textbf{Spatial Relationship Preservation:} Unlike 2D images where aggressive re-sampling or re-weighting primarily affects individual pixels, point clouds rely heavily on local neighborhood relationships. Naive application of 2D long-tailed techniques can disrupt these crucial spatial relationships, potentially degrading the learned geometric representations.

\textbf{Point Density Correlation:} In point clouds, class imbalance often correlates with point density variations. Tail classes are not only less frequent but may also be represented by fewer, sparser points, creating a compound challenge that traditional long-tailed methods do not address.

\textbf{Hierarchical Feature Learning:} Modern point-based networks like PointNet++ and ResPointNet++ use hierarchical feature learning with multiple levels of abstraction. It remains unclear which parts of these networks should be considered as "representation learning" versus "classification" components for decoupled training approaches.

These unique challenges motivate our work to develop ResPointNet-LT, which specifically addresses the adaptation of long-tailed learning techniques to hierarchical point-based networks while preserving the crucial spatial relationships that enable effective 3D understanding.



%% methodology
\section{Methodology}
\label{sec:methodology}
This section presents our ResPointNet-LT framework for addressing class imbalance in industrial point cloud semantic segmentation. The framework is designed to leverage the proven ResPointNet++ architecture while systematically tackling the long-tailed distribution problem through a novel two-stage decoupled training strategy. As illustrated in Figure~\ref{fig:framework}, our approach consists of three main components. First, we formulate the class imbalance problem and establish the industrial point cloud dataset characteristics (Section~\ref{sec:problem_formulation}). Second, we present the ResPointNet-LT architecture and two-stage decoupled training framework (Section~\ref{sec:respointnet_lt}), which separates representation learning from classifier training to achieve robust performance across all classes. Finally, we detail the training procedures, loss functions, and evaluation metrics specifically designed for assessing long-tailed performance (Section~\ref{sec:training_eval}).

\subsection{Problem Formulation and Dataset Characteristics}
\label{sec:problem_formulation}

\subsubsection{Long-Tailed Distribution in Industrial Point Clouds}
We formulate the class imbalance problem in industrial point cloud semantic segmentation as follows. Given an industrial point cloud dataset $\mathcal{D} = \{(P_i, Y_i)\}_{i=1}^N$ where $P_i \in \mathbb{R}^{M_i \times 3}$ represents the $i$-th point cloud with $M_i$ points and $Y_i \in \{1, 2, \ldots, C\}^{M_i}$ denotes the corresponding point-wise semantic labels for $C$ classes. Let $n_c$ be the total number of points belonging to class $c$ across the entire dataset, and $n = \sum_{c=1}^C n_c$ be the total number of labeled points.

The class frequency distribution $\pi_c = n_c/n$ exhibits a long-tailed pattern where classes can be categorized into three groups:
\begin{itemize}
    \item \textbf{Head classes} ($\mathcal{H}$): Classes with high frequency, typically structural elements like pipes, walls, and floors, where $\pi_c > \tau_h$ (e.g., $\tau_h = 0.1$).
    \item \textbf{Medium classes} ($\mathcal{M}$): Classes with moderate frequency, such as equipment supports and platforms, where $\tau_m < \pi_c \leq \tau_h$ (e.g., $\tau_m = 0.01$).
    \item \textbf{Tail classes} ($\mathcal{T}$): Classes with low frequency, typically critical components like pumps, valves, and sensors, where $\pi_c \leq \tau_m$.
\end{itemize}

The imbalance ratio $\rho = n_{\max}/n_{\min}$ where $n_{\max}$ and $n_{\min}$ are the frequencies of the most and least frequent classes, respectively. In industrial datasets, $\rho$ often exceeds 100:1, creating severe optimization challenges for standard training procedures.

\subsubsection{Industrial Point Cloud Dataset Characteristics}
Our evaluation focuses on large-scale industrial point cloud datasets that exhibit realistic long-tailed distributions. The primary dataset consists of point clouds from industrial facilities including petrochemical plants, power stations, and manufacturing facilities. Key characteristics include:

\begin{itemize}
    \item \textbf{Severe Class Imbalance}: Imbalance ratios ranging from 50:1 to 200:1, significantly exceeding typical computer vision datasets.
    \item \textbf{Critical Minority Components}: Tail classes often represent the most functionally important elements for industrial applications (pumps, valves, control systems).
    \item \textbf{Complex Geometric Structures}: Dense arrangements of pipes, equipment, and infrastructure with high levels of occlusion.
    \item \textbf{Varying Point Densities}: Different scanning resolutions and distances result in varying point densities across different components.
    \item \textbf{Large Scale}: Individual scenes containing millions of points spanning large industrial areas.
\end{itemize}

\subsubsection{Challenge Analysis}
The combination of severe class imbalance and 3D point cloud characteristics creates unique challenges:

1. \textbf{Representation Bias}: Standard training procedures lead to feature representations optimized for head classes, resulting in poor generalization to tail classes.
2. \textbf{Spatial Dependency}: Point clouds rely on local neighborhood relationships, making naive re-sampling strategies potentially harmful to geometric understanding.
3. \textbf{Compound Imbalance}: Tail classes are both less frequent and often represented by fewer, sparser points due to their smaller physical size or scanning artifacts.

\subsection{The Proposed ResPointNet-LT Framework}
\label{sec:respointnet_lt}

\subsubsection{Overview}
The ResPointNet-LT framework builds upon our established ResPointNet++ architecture and introduces a systematic two-stage decoupled training strategy to address class imbalance in industrial point cloud segmentation. The framework consists of two main components: (1) the ResPointNet++ backbone network that provides hierarchical feature learning capabilities, and (2) a novel two-stage training procedure that separates representation learning from classifier training. As illustrated in Figure~\ref{fig:respointnet_lt_framework}, our approach first trains the entire network on imbalanced data to learn robust geometric representations, then freezes the encoder and re-trains only the classifier components using balanced sampling strategies. This decoupling prevents the corruption of learned spatial representations while enabling the classifier to develop unbiased decision boundaries.

\subsubsection{ResPointNet++ Backbone Architecture}
\label{sec:respointnet_backbone}

Our framework builds upon the ResPointNet++ architecture, which extends the foundational PointNet++ with deep residual connections and enhanced local feature aggregation specifically designed for complex industrial scenes. The network employs a hierarchical encoder-decoder structure that captures multi-scale geometric features while maintaining spatial relationships crucial for 3D understanding.

\paragraph{Encoder Architecture}
The ResPointNet++ encoder consists of multiple set abstraction layers with residual connections. Each layer $i$ processes features from the previous layer through local aggregation and down-sampling:

\begin{equation}
F^{[i]} = \text{ResBlock}(\text{SA}(F^{[i-1]})) + \text{Shortcut}(F^{[i-1]})
\end{equation}

where $F^{[i]} \in \mathbb{R}^{N_i \times D_i}$ represents learned point embeddings for stage $i$, $F^{[0]}$ denotes the input point cloud $P \in \mathbb{R}^{N \times 3}$, SA represents the set abstraction operation, and ResBlock introduces residual learning to enable deeper architectures.

The set abstraction operation in each layer performs three key steps:
1. \textbf{Sampling}: Farthest point sampling (FPS) selects representative points from the input
2. \textbf{Grouping}: Ball query groups neighboring points around each representative point
3. \textbf{Feature Learning}: Local MLPs extract features from grouped points with max-pooling aggregation

\paragraph{Decoder Architecture}
The decoder uses feature propagation layers with skip connections to restore the original point resolution:

\begin{equation}
F_{\text{up}}^{[i]} = \text{FP}(F_{\text{up}}^{[i+1]}, F^{[i]}) = \text{MLP}(\text{Interp}(F_{\text{up}}^{[i+1]}) \oplus F^{[i]})
\end{equation}

where FP denotes feature propagation, Interp performs inverse distance weighted interpolation, and $\oplus$ represents feature concatenation. The final classification layer produces per-point predictions:

\begin{equation}
\hat{Y} = \text{MLP}_{\text{cls}}(F_{\text{up}}^{[0]}) \in \mathbb{R}^{N \times C}
\end{equation}

\paragraph{Architecture Components for Decoupled Training}
For our two-stage decoupled training strategy, we identify the network components as follows:
\begin{itemize}
    \item \textbf{Encoder Components}: All set abstraction layers, residual blocks, and feature extraction MLPs (parameters $\Theta_{\text{enc}}$)
    \item \textbf{Classifier Components}: Feature propagation layers, upsampling modules, and final classification MLPs (parameters $\Theta_{\text{cls}}$)
\end{itemize}

This decomposition enables us to separate representation learning (encoder) from decision boundary learning (classifier), which is crucial for addressing class imbalance without corrupting learned geometric features.

\subsubsection{Two-Stage Decoupled Training Strategy}
\label{sec:decoupled_training}

The core innovation of ResPointNet-LT lies in a systematic two-stage decoupled training approach that addresses class imbalance while preserving the crucial spatial relationships learned by the encoder. The training procedure is designed to leverage the full dataset during representation learning while enabling the classifier to develop unbiased decision boundaries.

\paragraph{Stage 1: Representation Learning}
In the first stage, we train the complete ResPointNet++ network end-to-end on the original imbalanced dataset using standard cross-entropy loss or class-balanced loss. The objective is to enable the encoder to learn rich geometric representations from all available data:

\begin{equation}
\mathcal{L}_{\text{stage1}} = \frac{1}{N} \sum_{i=1}^{N} \sum_{c=1}^{C} -y_{i,c} \log(\hat{y}_{i,c})
\end{equation}

where $y_{i,c}$ and $\hat{y}_{i,c}$ are the ground truth and predicted probabilities for point $i$ and class $c$. This stage typically runs for $E_1$ epochs (e.g., 200 epochs) until convergence.

The key insight is that despite the class imbalance, the encoder can still learn meaningful geometric features from all classes, as the spatial relationships and local geometric patterns are not fundamentally different between head and tail classes. The encoder learns to extract useful local and global features that capture the essential geometric structures in the point cloud.

\paragraph{Stage 2: Classifier Re-training}
In the second stage, we freeze the encoder parameters $\Theta_{\text{enc}}$ and re-train only the classifier components $\Theta_{\text{cls}}$ using a class-balanced sampling strategy:

\begin{equation}
\Theta_{\text{cls}}^* = \arg\min_{\Theta_{\text{cls}}} \mathcal{L}_{\text{balanced}}(\Theta_{\text{enc}}^*, \Theta_{\text{cls}})
\end{equation}

where $\Theta_{\text{enc}}^*$ are the frozen encoder parameters from Stage 1.

\textbf{Class-Balanced Sampling:} We implement a class-balanced sampling strategy that ensures each class has equal probability of appearing in training batches. For each training batch, we sample an equal number of points from each class, effectively creating balanced mini-batches for classifier training.

\textbf{Loss Function:} During Stage 2, we use standard cross-entropy loss since the input to the loss function is already balanced:

\begin{equation}
\mathcal{L}_{\text{balanced}} = \frac{1}{N_{\text{balanced}}} \sum_{i=1}^{N_{\text{balanced}}} \sum_{c=1}^{C} -y_{i,c} \log(\hat{y}_{i,c})
\end{equation}

This stage typically runs for $E_2$ epochs (e.g., 50 epochs) with a lower learning rate to fine-tune the decision boundaries.

\paragraph{Benefits of Decoupled Training}
The decoupled training strategy provides several advantages:

1. \textbf{Representation Preservation}: Freezing the encoder prevents corruption of learned geometric features during balanced training.
2. \textbf{Efficient Rebalancing}: Only the classifier needs to adapt to balanced data, reducing computational cost and training instability.
3. \textbf{Spatial Relationship Maintenance}: The encoder's understanding of local neighborhoods and spatial contexts remains intact.
4. \textbf{Flexible Loss Functions}: Different loss functions can be used in each stage to optimize different objectives.

\paragraph{Implementation Details}
\textbf{Parameter Initialization:} The classifier parameters in Stage 2 can be either kept from Stage 1 or re-initialized. Our experiments show that keeping Stage 1 parameters generally works better.

\textbf{Learning Rate Schedule:} We use a lower learning rate in Stage 2 (typically 1/10 of Stage 1 learning rate) to ensure stable fine-tuning.

\textbf{Data Augmentation:} Strong data augmentation (rotation, scaling, jittering) is applied in both stages, with additional augmentation for tail class samples in Stage 2.

\subsection{Training Procedures and Evaluation Metrics}
\label{sec:training_eval}

\subsubsection{Training Configuration}
\textbf{Two-Stage Training Schedule:}
\begin{itemize}
    \item \textbf{Stage 1}: Train complete ResPointNet++ for $E_1 = 200$ epochs with learning rate $\eta_1 = 0.001$, batch size 8, and standard data augmentation.
    \item \textbf{Stage 2}: Freeze encoder, re-train classifier for $E_2 = 50$ epochs with learning rate $\eta_2 = 0.0001$, using class-balanced sampling.
\end{itemize}

\textbf{Class-Balanced Sampling Strategy:} In Stage 2, each training batch contains equal number of points from each class. For classes with insufficient samples, we employ oversampling with augmentation to reach the required batch size.

\subsubsection{Loss Functions for Long-Tailed Learning}
We investigate multiple loss functions tailored for class imbalance:

\textbf{Standard Cross-Entropy Loss:}
\begin{equation}
\mathcal{L}_{CE} = -\frac{1}{N}\sum_{i=1}^N\sum_{c=1}^C y_{i,c} \log(\hat{y}_{i,c})
\end{equation}

\textbf{Class-Balanced Loss:} Inspired by~\cite{cui2019class}, we use effective number of samples to re-weight the loss:
\begin{equation}
\mathcal{L}_{CB} = -\frac{1}{N}\sum_{i=1}^N\sum_{c=1}^C \frac{1-\beta^{n_c}}{1-\beta} y_{i,c} \log(\hat{y}_{i,c})
\end{equation}
where $\beta \in [0,1)$ controls the re-weighting strength and $n_c$ is the number of samples in class $c$.

\textbf{Focal Loss:} To address the challenge of easy vs. hard examples:
\begin{equation}
\mathcal{L}_{Focal} = -\frac{1}{N}\sum_{i=1}^N\sum_{c=1}^C \alpha_c (1-\hat{y}_{i,c})^\gamma y_{i,c} \log(\hat{y}_{i,c})
\end{equation}
where $\alpha_c$ is the class weight and $\gamma$ is the focusing parameter.

\subsubsection{Evaluation Metrics for Long-Tailed Performance}
To comprehensively assess the performance of our ResPointNet-LT framework on imbalanced datasets, we employ a comprehensive set of evaluation metrics that provide insights into both overall performance and class-specific behavior, with particular emphasis on tail class performance.

\textbf{Standard Segmentation Metrics:}

\textbf{Per-Class IoU:} Intersection over Union for each class $c$:
\begin{equation}
\label{eq:iou}
\text{IoU}_c = \frac{TP_c}{TP_c + FP_c + FN_c}
\end{equation}

\textbf{Mean IoU (mIoU):} Macro-averaged IoU across all classes:
\begin{equation}
\label{eq:mIoU}
\text{mIoU} = \frac{1}{C}\sum_{c=1}^{C} \text{IoU}_c
\end{equation}

\textbf{Overall Accuracy (OA):} Percentage of correctly classified points:
\begin{equation}
\label{eq:accuracy}
\text{OA} = \frac{\sum_{c=1}^{C} TP_c}{\sum_{c=1}^{C} (TP_c + FP_c + FN_c)}
\end{equation}

\textbf{Long-Tailed Specific Metrics:}

\textbf{Head-mIoU:} Mean IoU over head classes only:
\begin{equation}
\text{Head-mIoU} = \frac{1}{|\mathcal{H}|}\sum_{c \in \mathcal{H}} \text{IoU}_c
\end{equation}

\textbf{Tail-mIoU:} Mean IoU over tail classes only (our key metric):
\begin{equation}
\text{Tail-mIoU} = \frac{1}{|\mathcal{T}|}\sum_{c \in \mathcal{T}} \text{IoU}_c
\end{equation}

\textbf{Balanced Accuracy:} Macro-averaged per-class accuracy:
\begin{equation}
\text{Balanced Accuracy} = \frac{1}{C}\sum_{c=1}^{C} \frac{TP_c}{TP_c + FN_c}
\end{equation}

\textbf{Performance Gap:} Difference between head and tail class performance:
\begin{equation}
\text{Performance Gap} = \text{Head-mIoU} - \text{Tail-mIoU}
\end{equation}

The Tail-mIoU is our primary metric for evaluating long-tailed performance, as it directly measures the model's ability to recognize minority classes that are often most critical in industrial applications. The Performance Gap metric quantifies the bias toward head classes, with smaller values indicating better balanced performance.

%% experimental results
\section{Experimental Results}
\label{sec:results}

This section presents the comprehensive experimental evaluation of the proposed ResPointNet-LT framework on industrial point cloud datasets with severe class imbalance. We first detail the experimental setup, datasets, and baseline methods (\autoref{sec:exp_config}). We then compare the performance of ResPointNet-LT against strong baselines including our previous ResPointNet++ work (\autoref{sec:main_results}). Extensive ablation studies analyze the contribution of each component in our framework (\autoref{sec:ablation_studies}). Finally, we provide qualitative analysis and visualization results that demonstrate the improved performance on tail classes (\autoref{sec:qualitative_analysis}). \autoref{tab:main_results} summarizes the key performance comparisons.

\subsection{Experimental Configuration}
\label{sec:exp_config}

\subsubsection{Dataset Description}
Our primary evaluation is conducted on a large-scale industrial point cloud dataset collected from petrochemical and power generation facilities. The dataset contains XX point clouds with a total of XX million points across XX classes, exhibiting severe long-tailed distribution with imbalance ratios ranging from 50:1 to 200:1.

\textbf{Dataset Statistics:}
\begin{itemize}
    \item \textbf{Head classes (5 classes):} pipe, wall, floor, structure, platform - comprising 85\% of all points
    \item \textbf{Medium classes (7 classes):} tank, support, stairs, railing, etc. - comprising 13\% of all points  
    \item \textbf{Tail classes (8 classes):} pump, valve, sensor, control\_panel, instrument, etc. - comprising 2\% of all points
    \item \textbf{Training/Validation/Test split:} 60\%/20\%/20\%
    \item \textbf{Point density:} 500-2000 points per m³
    \item \textbf{Scene complexity:} High occlusion, dense pipe networks, multi-level structures
\end{itemize}

\subsubsection{Baseline Methods}
We compare ResPointNet-LT against several strong baselines:
\begin{itemize}
    \item \textbf{ResPointNet++ (Baseline):} Our previous work without long-tailed learning techniques
    \item \textbf{PointNet++:} The foundational hierarchical point network
    \item \textbf{RandLA-Net:} Efficient large-scale point cloud segmentation network
    \item \textbf{ResPointNet++ + CB-Loss:} ResPointNet++ with class-balanced loss
    \item \textbf{ResPointNet++ + Focal Loss:} ResPointNet++ with focal loss
    \item \textbf{ResPointNet++ + Oversampling:} ResPointNet++ with simple oversampling of tail classes
\end{itemize}

\subsubsection{Implementation Details}
All experiments are conducted using PyTorch on NVIDIA V100 GPUs. Training configurations follow the two-stage protocol described in Section~\ref{sec:training_eval}. We report results averaged over 3 independent runs with different random seeds.

\subsection{Main Results}
\label{sec:main_results}

\textbf{Overall Performance:} Table~\ref{tab:main_results} shows that ResPointNet-LT significantly outperforms all baseline methods. Our approach achieves XX.X\% mIoU, representing a XX.X\% improvement over the ResPointNet++ baseline, while maintaining competitive overall accuracy.

\textbf{Tail Class Performance:} The most significant improvement is observed in tail class performance. ResPointNet-LT achieves XX.X\% Tail-mIoU compared to XX.X\% for ResPointNet++ baseline - a remarkable improvement of XX.X percentage points. Individual tail class improvements include:
\begin{itemize}
    \item Pump: XX.X\% → XX.X\% IoU (+XX.X\%)
    \item Valve: XX.X\% → XX.X\% IoU (+XX.X\%)  
    \item Sensor: XX.X\% → XX.X\% IoU (+XX.X\%)
    \item Control Panel: XX.X\% → XX.X\% IoU (+XX.X\%)
\end{itemize}

\textbf{Performance Gap Reduction:} The performance gap between head and tail classes is reduced from XX.X\% to XX.X\%, demonstrating more balanced performance across all classes.

\subsection{Ablation Studies}
\label{sec:ablation_studies}

\subsubsection{Component Analysis}
We conduct comprehensive ablation studies to analyze the contribution of each component:

\textbf{Training Strategy Comparison:}
\begin{itemize}
    \item Single-stage training: XX.X\% mIoU, XX.X\% Tail-mIoU
    \item Two-stage decoupled training: XX.X\% mIoU, XX.X\% Tail-mIoU (+XX.X\%)
\end{itemize}

\textbf{Loss Function Analysis:}
\begin{itemize}
    \item Cross-entropy loss: XX.X\% Tail-mIoU
    \item Class-balanced loss: XX.X\% Tail-mIoU (+XX.X\%)
    \item Focal loss: XX.X\% Tail-mIoU (+XX.X\%)
    \item Combined approach: XX.X\% Tail-mIoU (+XX.X\%)
\end{itemize}

\subsubsection{Hyperparameter Sensitivity}
Analysis of key hyperparameters including stage training epochs, learning rates, and class-balanced sampling strategies.

\subsection{Qualitative Analysis and Visualization}
\label{sec:qualitative_analysis}

Figure~\ref{fig:qualitative_results} presents side-by-side visualizations comparing ResPointNet++ baseline with our ResPointNet-LT on representative industrial scenes. Key observations:

\begin{itemize}
    \item \textbf{Pump Identification:} ResPointNet-LT correctly identifies small pumps that were previously misclassified as pipes or structural elements
    \item \textbf{Valve Recognition:} Significant improvement in valve segmentation accuracy, especially for small control valves
    \item \textbf{Sensor Detection:} Better detection of various sensor types mounted on equipment
    \item \textbf{Preserved Head Class Performance:} No degradation in the segmentation quality of dominant structural elements
\end{itemize}

These qualitative results demonstrate that ResPointNet-LT successfully addresses the class imbalance problem while maintaining the overall segmentation quality essential for industrial applications.

%% conclusion
\section{Conclusion}
\label{sec:conclusion}

This paper presents ResPointNet-LT, a novel two-stage decoupled training framework specifically designed to address the critical challenge of class imbalance in industrial point cloud semantic segmentation. Building upon our established ResPointNet++ architecture, we systematically tackle the long-tailed distribution problem that severely impacts the recognition of critical minority components in real-world industrial applications.

Our key contributions include: (1) a comprehensive analysis and formulation of the class imbalance problem in industrial point cloud datasets, highlighting the unique challenges compared to traditional computer vision datasets; (2) a pioneering two-stage decoupled training strategy that separates representation learning from classifier training, enabling effective handling of class imbalance while preserving crucial spatial relationships; (3) adaptation of proven long-tailed learning techniques from 2D computer vision to the 3D point cloud domain, addressing specific challenges related to spatial dependencies and hierarchical feature learning; (4) extensive experimental validation demonstrating substantial improvements in tail class performance while maintaining head class accuracy.

The experimental results confirm the effectiveness of our approach, with ResPointNet-LT achieving significant improvements in Tail-mIoU (XX.X percentage points improvement) and substantial reductions in performance gap between head and tail classes. The qualitative analysis further demonstrates the practical value of our method for real-world industrial applications, where accurate identification of minority components such as pumps, valves, and sensors is crucial for facility management and maintenance.

\textbf{Limitations and Future Work:} While ResPointNet-LT demonstrates significant improvements, several limitations remain. The current approach requires manual definition of head/tail class thresholds, and the optimal training schedule may vary across different industrial domains. Future research directions include: (1) developing adaptive threshold selection mechanisms for class categorization; (2) investigating dynamic balancing strategies that adjust throughout training; (3) extending the approach to other point-based architectures beyond ResPointNet++; (4) exploring the integration of semantic knowledge and geometric constraints to further improve tail class recognition; (5) conducting large-scale evaluations across diverse industrial domains to assess generalizability.

In conclusion, ResPointNet-LT represents a significant step toward practical deployment of deep learning techniques in real-world industrial point cloud analysis, providing a robust solution to the ubiquitous class imbalance challenge while maintaining the high-quality geometric understanding essential for complex 3D industrial scenes.

%%%%%%%%%%%%%%%%%%%%%%%%%%%%%%%%%%%%%%%%%%
\stopmainlinenumbers

\section*{Acknowledgments}
This work was supported by the China Postdoctoral Science Foundation (No.~2023M740761), the National Natural Science Foundation of China (No.~42001198), the Natural Science Foundation of Hunan Province, China (No.~2023JJ40098), and the GDAS’ Project of Science and Technology Development (No.~2022GDASZH-2022010202, 2022GDASZH-2022020402-01, 2022GDASZH-2022010111).

\section*{Data Availability Statement}
The authors confirm that the data supporting the findings of this study are available within the article.

\section*{Declaration of Competing Interest}
The authors declare that they have no known competing financial interests or personal relationships that could have appeared to influence the work reported in this paper.

%% references
\bibliographystyle{elsarticle-num}
\bibliography{references}

\end{document}

\endinput